% TOP_PROBLEM: {"id": 30006691, "title": "The Quantum PCP Conjecture", "category_id": 15, "category": {"id": 15, "name": "computer_science", "display_name": "Computer Science", "description": "Computational complexity, algorithms, and theoretical CS.", "slug": "computer-science", "order_index": 15, "created_at": "2026-07-31T15:26:25.671Z"}, "status": "open"}
% FIRST_POSTED: 2026-09-27
% !TeX program = lualatex
% Compile twice: lualatex open_problems_500.tex
% All references and prose are embedded; no BibTeX database or external inputs.
\documentclass[11pt,a4paper]{article}
\usepackage[margin=24mm,headheight=14pt]{geometry}
\usepackage{fontspec}
\setmainfont{Latin Modern Roman}
\setsansfont{Latin Modern Sans}
\setmonofont{Latin Modern Mono}
\usepackage{amsmath,amssymb,mathtools}
\usepackage{unicode-math}
\setmathfont{Latin Modern Math}
\usepackage{microtype}
\usepackage{enumitem,xurl,needspace}
\usepackage[dvipsnames]{xcolor}
\usepackage{hyperref}
\hypersetup{unicode=true,colorlinks=true,linkcolor=MidnightBlue,urlcolor=MidnightBlue,
 pdftitle={Research notebook: Problems 40--49},pdfauthor={Source-annotated compilation},bookmarksnumbered=true}
\usepackage{bookmark}
\usepackage{tocloft}
\setlength{\cftsecnumwidth}{3em}
\cftsetpnumwidth{2.5em}
\cftsetrmarg{4em}
\usepackage{titlesec}
\titleformat{\section}{\Large\bfseries\raggedright}{\thesection}{0.7em}{}
\usepackage{fancyhdr}
\pagestyle{fancy}\fancyhf{}
\fancyhead[L]{\small\sffamily Mathematical problem statements}
\fancyhead[R]{\small Source edition 22}
\fancyfoot[C]{\thepage}
\setlength{\parindent}{0pt}
\setlength{\parskip}{5pt plus 1pt}
\setlength{\emergencystretch}{3em}
\setcounter{tocdepth}{1}
\setcounter{secnumdepth}{1}
\setlist[itemize]{leftmargin=3em,itemsep=5pt,topsep=4pt,parsep=0pt}
\allowdisplaybreaks
\newcommand{\N}{\mathbb{N}}
\newcommand{\Z}{\mathbb{Z}}
\newcommand{\Q}{\mathbb{Q}}
\newcommand{\R}{\mathbb{R}}
\newcommand{\C}{\mathbb{C}}
\newcommand{\F}{\mathbb{F}}
\newcommand{\A}{\mathbb{A}}
\newcommand{\PP}{\mathbb P}
\newcommand{\E}{\mathbb{E}}
\newcommand{\eps}{\varepsilon}
\newcommand{\dd}{\,\mathrm d}
\newcommand{\sref}[2]{\hyperlink{source:#1:#2}{[S#2]}}
\newcommand{\eref}[2]{\hyperlink{extra:#1:#2}{[E#2]}}
% Turn the next switch off for a shorter reading copy. The quotations preserve
% every exactTarget field, including qualifications omitted from a short summary.
\newif\ifshowcatalogue
\showcataloguetrue
\newcommand{\cataloguescope}[1]{\ifshowcatalogue
 \paragraph{Exact scope in the supplied catalogue.}
 \begin{quote}\small #1\end{quote}\fi}

% Research additions dated 27 September 2026; compile twice with LuaLaTeX.
\numberwithin{equation}{section}
\begin{document}
\setcounter{section}{0}
% ================================================================
% problemId: 30006691
% Canonical problem ID: 30006691
% exactTarget SHA-256: 51038dfdfa4c11daba96c361d6d4ec63f18fc276d233b3fb17fa991001caad69
\clearpage
\section*{\texorpdfstring{The Quantum PCP Conjecture}{The Quantum PCP Conjecture}}\label{30006691}
\subsection{Definitions and mathematical statement}
A $k$-local Hamiltonian on $n$ finite-dimensional sites is $H=\sum_{j=1}^m h_j$, each term acting on at most a constant $k$ sites. Fix a normalization such as $0\le h_j\le I$ and define ground energy $\lambda_{\min}(H)$. The quantum PCP assertion asks for constants $k$ and $0\le a<b\le1$ for which
\[
 \lambda_{\min}(H)\le am\quad\text{versus}\quad\lambda_{\min}(H)\ge bm
\]
is QMA-hard, under the standard efficient reductions. The separation $(b-a)m$ is a constant fraction of the number of normalized terms, not an inverse-polynomial absolute gap. The exact Hamiltonian encoding and promise convention are fixed by the cited formulation.
\par\smallskip\textit{Formulation and concept references:} \sref{30006691}{1}.
\cataloguescope{Prove or disprove that approximating the ground-state energy of a local Hamiltonian to within a constant additive fraction of its number of local terms is QMA-hard.}
\subsection{Short English statement}
Is estimating a normalized local quantum system's energy to constant relative precision already as hard as general quantum-witness verification?
% BEGIN RESEARCH ADDITION 41
\subsection{Research attempt}
\paragraph{Current frontier.}
\textit{Literature check: 27 September 2026.}
The target is the constant-locality, constant \emph{normalized energy}
promise of Aharonov--Arad--Vidick, Section 1.2, Conjecture 1.3.
Neither an absolute spectral gap nor a difficult-to-prepare ground state
alone supplies the required QMA-hard reduction. \sref{30006691}{1}.
The NLTS theorem of Anshu--Breuckmann--Nirkhe supplies Hamiltonians whose
low-energy states require nonconstant circuit depth, but does not itself
prove this Hamiltonian quantum PCP statement. \eref{30006691}{1}.
A 2025 manuscript by Bergamaschi--Metger--Vidick--Zhang gives derandomized
amplification and a locality--gap tradeoff. Its introduction explicitly
identifies locality reduction as the missing part of a full PCP route;
its streaming construction does not have constant-locality terms.
\eref{30006691}{2}. These results are kept separate from the construction below.

\paragraph{Attack route.}
Duplicating constraints multiplies both energy and term count and therefore
does not amplify normalized energy. Applying a nonlinear spectral function
can amplify small eigenvalues, but may destroy locality. A third route is
to encode into a quantum code and test locally; this additionally needs a
soundness theorem that rules out all entangled low-energy cheating states.
We investigate the spectral route and try to repair the positivity and
noncommutativity problems before confronting locality. No efficient
locality-reduction theorem is assumed.

\paragraph{Attempt: positive word amplification without commutativity.}
Let $H=\sum_{i=1}^m h_i$, $0\le h_i\le I$, and put
$B=H/m$, $C_i=I-h_i$. For a word $w=(i_1,\ldots,i_t)$ define
\[
 W_w=C_{i_1}\cdots C_{i_t},\qquad
 g_w=\frac{2I-W_w-W_w^*}{4}.
 \tag{41.1}
\]
Each $C_i$ is a contraction, so $\|W_w\|\le1$.
For every unit vector $v$,
$-1\le\operatorname{Re}\langle v,W_wv\rangle\le1$;
therefore $0\le g_w\le I$. Moreover $g_w$ is supported on the union of
at most $t$ original supports, hence is at most $kt$-local.
This proof does not assert that the product of positive matrices is
positive or even Hermitian.

Averaging over all ordered words gives the exact operator identity
\[
 G_t:=\frac1{m^t}\sum_{w\in[m]^t}g_w
   =\frac12\bigl[I-(I-B)^t\bigr].
 \tag{41.2}
\]
Indeed expansion of $(\sum_i C_i)^t$ includes all ordered products in
exactly their original order; commutation is never used. The adjoint sum
is identical because reversing words permutes the word set. By the
spectral theorem and monotonicity on $[0,1]$,
\[
 \lambda_{\min}(G_t)=f_t(e),\qquad
 e=\lambda_{\min}(H)/m,\qquad
 f_t(u)=\tfrac12[1-(1-u)^t].
 \tag{41.3}
\]
The unnormalized Hamiltonian $\sum_wg_w$ has $m^t$ terms, so (41.3)
really is its energy divided by its term count, as required by the
original normalization.

Suppose a hard input family had normalized thresholds $0\le a\le b/100$
and $0<b\le1$. With $t=\lceil1/b\rceil$,
\[
 f_t(a)\le\frac{ta}{2}\le\frac1{100},\qquad
 f_t(b)\ge\frac{1-e^{-1}}2.
 \tag{41.4}
\]
Here $t\le2/b$, $1-(1-a)^t\le ta$, and
$(1-b)^t\le e^{-tb}\le e^{-1}$. Thus the spectral transformation has
an absolute normalized separation exceeding $0.30$ in this parameter
regime. This is an exact mathematical amplification lemma, not yet an
efficient reduction. The small-completeness-energy condition is stated
rather than replaced by perfect completeness, which would be an additional
assumption on a general QMA verification problem.

For arbitrary $a<b$, the output separation is instead
\[
 f_t(b)-f_t(a)=\tfrac12[(1-a)^t-(1-b)^t].            \tag{41.5}
\]
If $a>0$ is fixed this tends to zero as $t\to\infty$. Thus ``take a large
power'' is not by itself a gap-amplification theorem for every promised
input. Threshold placement matters.

\paragraph{Attempt: can locality be recovered for free?}
We next test the strongest naive hope: perhaps the operator in (41.2)
always admits some other bounded-locality decomposition on the same
qubits. This is false, already for commuting one-local terms.
Take $n\ge t$ qubits and
\[
 h_i=|1\rangle\langle1|_i,\qquad B=\frac1n\sum_{i=1}^n h_i.
\]
On a computational basis vector $x\in\{0,1\}^n$, $G_t$ has diagonal
entry $f_t(|x|/n)$. In the unique multilinear polynomial representing
this function on the Boolean cube, the coefficient of
$x_{i_1}\cdots x_{i_t}$, for distinct indices, is
\[
 c_t=\frac{(-1)^{t+1}t!}{2n^t}\ne0.               \tag{41.6}
\]
To see this, expand $f_t(u)$ as a degree-$t$ polynomial. Only its leading
term can produce a monomial with $t$ distinct indices in
$(\sum x_i)^t$; there are $t!$ contributing orders.

The diagonal entries of an operator supported on at most $k$ qubits
depend on at most $k$ Boolean coordinates and have multilinear degree at
most $k$. A sum of such operators has the same degree bound. Taking
$t=k+1$ in (41.6) therefore proves that no sum of $k$-local operators,
positive or otherwise, equals $G_t$ on the same Hilbert space.
Off-diagonal entries of a proposed decomposition cannot change this
diagonal obstruction.

Equivalently the Pauli coefficient of
$Z_{i_1}\cdots Z_{i_t}$ is
\[
 -\frac{t!}{2^{t+1}n^t}.
 \tag{41.7}
\]
Consequently any $k$-local $K$ on the same qubits, for $k<t$, satisfies
\[
 \|G_t-K\|\ge \frac{t!}{2^{t+1}n^t}.
\]
This follows from
$2^{-n}|\operatorname{Tr}(Z_{i_1}\cdots Z_{i_t}(G_t-K))|
\le\|G_t-K\|$. The lower bound is not asserted to stay constant as the
parameters grow. It is an obstruction to \emph{exact} locality recovery,
not a theorem ruling out useful approximate gadget constructions with
additional qubits.

\paragraph{Stress test: size, locality, and encodings.}
There are two independent efficiency failures. When $b$ is inverse
polynomial, $t$ is polynomial rather than constant, and $kt$ need not be
bounded. Also $m^t$ is generally exponential. A short symbolic expression
for (41.2) is not the polynomial-length list of bounded-locality matrices
required by the target. Iterating the full-word construction at constant
$t$ does not cure the size problem: after $r$ rounds the number of terms
is $m^{t^r}$. This is why the derandomized construction in the recent
literature supplies genuine additional work. \eref{30006691}{2}.

The finite tests accompanying this entry use the noncommuting rational
projectors $|0\rangle\langle0|$ and $|+\rangle\langle+|$. They verify
(41.2) exactly for short words and check positivity of every summand.
For three commuting one-qubit terms they verify the weight-three
coefficient $-1/72$ when $t=n=3$. These are sanity checks of identities
already proved above, not hardness experiments.

An ancilla construction could evade (41.6), because the effective energy
on a subspace or after minimizing over extra qubits is not a literal
$k$-local operator on the original register. It would, however, need all
of the following in one deterministic polynomial-time reduction:
constant final locality, a polynomial explicit term list with norms at
most one, completeness below a fixed threshold, and soundness above a
strictly larger fixed threshold for \emph{every} enlarged-system state.
An error estimate in an unnormalized low-energy effective Hamiltonian is
insufficient if dividing by the gadget's total term count erases the gap.
No such reduction is proved here.

Finally, QMA-hardness itself is not a separation from BQP. If a separate
argument put the promised approximation problem in BQP, quantum PCP would
imply QMA$=$BQP rather than a formal contradiction. This is the relevant
dependency on UnsolvedMath problem 30006695. Conversely, the NLTS theorem concerns preparation
depth and cannot simply be substituted for a decision-class lower bound.
\sref{30006691}{1}, \eref{30006691}{1}.

\paragraph{Outcome.}
\textbf{Outcome: Exploratory approach with unresolved gap.}
The positive-word identity provides a fully verified spectral
amplification calculation for noncommuting constraints. The Boolean-degree
test proves that exact locality recovery on the unchanged register fails
even in a simple commuting example. Neither lemma is presented as a new
historical result or as quantum PCP. The unresolved step is a polynomial,
constant-locality realization with a surviving normalized promise gap;
the explicit construction here has neither the necessary term-count bound
nor the necessary locality bound.
% END RESEARCH ADDITION 41
\subsection{Sources}
\begin{itemize}\raggedright
\item[\textnormal{[S1]}]\hypertarget{source:30006691:1}{} Aharonov, Arad, and Vidick, The Quantum PCP Conjecture, Section 1.2, Conjecture 1.3 (2013).
\url{https://arxiv.org/html/1309.7495v1}.
{\small\textit{Catalogue formulation source.}}
\item[\textnormal{[S2]}]\hypertarget{source:30006691:2}{} Private PCPs from Product Expansion.
\url{https://eccc.weizmann.ac.il/report/2026/150/}.
{\small\textit{Catalogue research/status reference; not a proof endorsement.}}
% BEGIN RESEARCH SOURCES 41
\item[\textnormal{[E1]}]\hypertarget{extra:30006691:1}{} Anurag Anshu, Nikolas P. Breuckmann,
and Chinmay Nirkhe, \emph{NLTS Hamiltonians from good quantum codes},
arXiv:2206.13228, STOC 2023, pp. 1090--1096; main NLTS theorem.
\url{https://arxiv.org/abs/2206.13228}.
{\small\textit{Established partial result, not a quantum PCP theorem;
consulted 27 September 2026.}}
\item[\textnormal{[E2]}]\hypertarget{extra:30006691:2}{} Thiago Bergamaschi, Tony Metger,
Thomas Vidick, and Tina Zhang, \emph{Derandomised tensor product gap
amplification for quantum Hamiltonians}, arXiv:2510.01333v1,
1 October 2025. Section 1, Theorems 1.2--1.4, Definition 1.6,
and the locality-increase discussion.
\url{https://arxiv.org/html/2510.01333v1}.
{\small\textit{Research manuscript; the notebook uses its stated scope and
obstructions, not an unverified assertion of quantum PCP.
Consulted 27 September 2026.}}
% END RESEARCH SOURCES 41
\end{itemize}

\end{document}
